\documentclass[12pt]{article}
\usepackage[utf8]{inputenc}
\usepackage[T1]{fontenc}
\usepackage[margin=1in]{geometry}
\usepackage{setspace}
\usepackage{xcolor}
\usepackage[colorlinks=true, urlcolor=blue!60!black, linkcolor=blue!60!black]{hyperref}
\setlength{\parindent}{0em}
\setlength{\parskip}{0.6em}
\renewcommand{\baselinestretch}{1.0}
\begin{document}
\begin{center}
{\LARGE \textit{Wolford v. Lopez}, 609 U.S. 185 (2026)\par}
\vspace{0.5em}
{\large Case Brief --- Two-Page Edition\par}
\end{center}

\section{Memory Jogger}
Hawaii's ``vampire rule,'' barring licensed carry onto property open to the public without the owner's express consent, violates the Second Amendment (\href{https://www.supremecourt.gov/opinions/25pdf/24-1046_nmio.pdf}{Slip Op.}).

\section{Detailed Case Facts}
After \textit{Bruen} (2022), Hawaii's Act 52 tightened licensing, banned guns in many sensitive places, and in \S{}134-9.5 forbade licensed carriers to bring firearms onto ``private property of another'' without express authorization by sign or word. This reversed the common-law implied license to enter property open to the public. Three Maui permit holders and the Hawaii Firearms Coalition sued Attorney General Anne E. Lopez.

\section{Procedural History}
The district court enjoined the rule for property open to the public; the Ninth Circuit reversed, finding a historical tradition of consent defaults. The Supreme Court reversed on June 25, 2026.

\section{Judicial Votes}
Reversed, 6--3: Alito wrote for Roberts, Thomas, Gorsuch, Kavanaugh, and Barrett; Barrett concurred; Kagan and Jackson (with Sotomayor) dissented.

\section{Holding}
A state may not require express owner consent before licensed carriers bring handguns onto private property open to the public. The rule burdens conduct within the Second Amendment's text and lacks a relevantly similar historical analogue; owners keep their own right to exclude.

\section{Analysis of Opinions}
Alito placed property-law history at \textit{Bruen} step two and showed that a carrier running ordinary errands could become ``a criminal at least six times over.'' The right has one national meaning; colonial laws were anti-poaching measures, and an 1865 Louisiana Black Code was a ``tainted artifact.'' Barrett stressed that the Amendment binds state property law and that founding-era laws support only focused defaults for places prone to abuse. Kagan would uphold at step two under \textit{Rahimi}; Jackson called the case ``about property rights, not gun rights'' and accused the majority of creating a ``one-way ratchet.''

\section{Examples: Future Applications}
\textbf{Same-Side:} a signage-only opt-in rule for businesses is invalid. \textbf{Opposite-Side:} a business's own ``No Firearms'' sign enforced through trespass law is valid. \textbf{Fence-Sitter:} an opt-in default limited to nightclubs and cannabis dispensaries after shootings.

\section{Critique}
Supporters such as Joseph Greenlee praise the decision for stopping post-\textit{Bruen} evasion. Maureen Brady and the \textit{Harvard Law Review} see a level-of-generality problem and an unguided test for discarding tainted historical laws.

\section{Key Quotations}
\begin{quote}
``This regime hobbles what the Second Amendment protects: the right of Americans to carry arms for self-defense as they go about their daily lives. We hold that the law is unconstitutional.''
\end{quote}
\noindent --- Justice Alito, Opinion of the Court.

\end{document}
